\chapter{Conclusiones y trabajo futuro}
\label{ch:conclusiones}

\section{Conclusiones}
\label{sec:concl-conclusiones}

Este trabajo se propuso predecir el desplazamiento diario de la gaviota sombría a
partir de su rastro GPS, mediante un enfoque híbrido que combina modelos
probabilísticos y aprendizaje automático, y comparar de forma sistemática las
soluciones desarrolladas bajo una métrica común y frente a baselines explícitas. Ese
objetivo general se ha alcanzado: el sistema encadena las distintas piezas en un flujo
coherente y las evalúa todas con el mismo criterio, de modo que su comparación es
directa.

Los cinco objetivos específicos planteados en la introducción
(\ref{sec:intro-objetivos}) se han cumplido. La preparación de los datos transformó los
registros GPS en bruto en una secuencia diaria homogénea por individuo, con los huecos
del seguimiento representados de forma explícita, que ha servido de base común a todos
los modelos (capítulo \ref{ch:o1}). Sobre ella, la cadena de Markov visible estableció
una baseline interpretable y fijó el criterio de evaluación, el log-loss, junto con la
referencia de persistencia frente a la que se contrasta el resto (capítulo
\ref{ch:o2}). El modelo oculto de Markov recuperó sin supervisión los dos regímenes de
comportamiento, estacionario y de migración, y se validó por su coherencia con la
fenología conocida de la especie (capítulo \ref{ch:o3}). Los modelos de aprendizaje
supervisado predijeron el desplazamiento empleando únicamente información disponible
antes del instante predicho e incorporando el régimen inferido, y entregaron una
predicción acompañada de una banda de incertidumbre bien calibrada (capítulo
\ref{ch:o4}). Por último, la herramienta de cartografía interactiva permite visualizar y
comparar sobre el mapa las predicciones de los distintos modelos (capítulo
\ref{ch:o5}).

De la comparación sistemática se extrae la conclusión central del trabajo. Los modelos
de aprendizaje supervisado se construyeron con el propósito de superar a la baseline de
Markov, y lo logran: la mejoran de forma consistente bajo el criterio común de
evaluación. Sobre esa base, reformular el problema como una regresión de cuantiles
permitió ir un paso más allá, al recuperar la precisión en la predicción del destino
diario que la primera formulación discretizada había cedido y, sobre todo, al entregar
una banda de incertidumbre bien calibrada, una salida directamente cartografiable que la
herramienta interactiva lleva al mapa. El enfoque híbrido funciona, por tanto, como un
sistema integrado y cumple lo que se propuso.

La evaluación delimita además, con datos, hasta dónde puede llegar la predicción a esta
escala: el horizonte de un día con información puramente local marca un techo
estructural, porque la posición y el contexto de una jornada no contienen el predictor de
un evento migratorio (el viento, la historia de varios días, el destino). Lejos de ser
una limitación del diseño, caracterizar ese techo orienta el trabajo futuro y sitúa el
valor del enfoque donde de verdad aporta: en la incertidumbre calibrada y en distinguir los días en que el ave
está migrando, que son los pocos en los que predecir su desplazamiento aporta
información real, frente a los muchos días en que apenas se mueve.

El trabajo deja, además, dos lecciones metodológicas que trascienden este caso concreto
y cuya proyección se analiza en el capítulo \ref{ch:impacto}: el diseño causal de las
variables, que exige que toda variable predictiva sea calculable sin observar el futuro,
y el uso de la coherencia biológica como criterio de validación, que contrasta cada
resultado contra la fenología de la especie antes de darlo por bueno.

\section{Trabajo futuro}
\label{sec:concl-futuro}

\todo[inline]{Trabajo futuro: modelos de secuencia, viento causal, multi-paso + reabribles.}
